
\def\PUDcodigo{EME120}

\if\COMPILAR0
   \def\PUDcodacad{04507.49}
\else
   \def\PUDcodacad{04506.48}
\fi

\def\PUDdisciplina{Gestão e Controle da Qualidade}

\def\PUDsemestre{S10}

\def\PUDhoras{80}

%NOVOS ---------------------------------------------------
\def\PUDhorasTeoria{80}
\def\PUDhorasPratica{0}

\if\COMPILAR0
   \def\PUDinterECA{ \hyperref[PUD:EME127]{Gestão da Manutenção Industrial (04507.50)} }
\else
   \def\PUDinterECA{ \hyperref[PUD:EME127]{Gestão da Manutenção Industrial (04506.49)} }
\fi

\def\PUDatuaisECA{---}

\def\PUDrecursos{Projetor multimídia; Quadro branco e pincel.}

%----------------------------------------------------------

\def\PUDcreditos{4}

\def\PUDprerequisito{---}

\def\PUDpreacad{---}

\def\PUDobjetivos{Entender os conceitos que envolvem a gestão da qualidade total na indústria;  Conhecer as práticas e ferramentas de um sistema da qualidade na indústria;  Conhecer a metodologia de implantação de um sistema de gestão da qualidade na indústria.}

\def\PUDementa{Gestão da Qualidade;  Metodologia para Programas de Melhoria da Qualidade;  Técnicas para Análise de Processos;  Certificação ISO;  Condições de Certificação;  Auditorias da Qualidade;  Plano de Ação para Certificação.}

\def\PUDconteudo{ & História e evolução da qualidade.  Principais conceitos em qualidade.  5S. 
\\ 
 & Planejamento e gestão da qualidade.  Satisfação do cliente. 
\\ 
 & Contribuições para a gestão da qualidade.  Qualidade ao estilo Japonês. 
\\ 
 & Processos.  Diagrama causa e efeito. 
\\ 
 & Métodos de controle dos processos.  Ciclo PDCA. 
\\ 
 & Técnicas para análise dos processos.  Cadeia de valor.  Fluxogramas. 
\\ 
 & Especificação e normas.  Conformidade.  Padronização. 
\\ 
 & ISO 9001:2000.  Qualidade na ótica do cliente. 
\\ 
 & Auditoria da qualidade.  Relatório da qualidade. 
\\ 
 & Gestão da qualidade em prestação de serviços. 
\\ 
 & Práticas do controle da qualidade.  Planos de ação para melhorias.  Ações preventivas e corretivas. 
\\ 
 & Benchmarking.  Kaizen. 
\\ 
 & Liderança.  Comprometimento.  Treinamento.  Cultura da qualidade. 
\\ 
 & Análise e solução de problemas.  Análise de Pareto. 
\\ 
 & Responsabilidade social.  Sustentabilidade. 
\\ 
 & Seminário de implantação da norma ISO 9001.  Seminário de certificação ISO 9001. }

\def\PUDmetodo{Aulas expositórias que podem ser teóricas e/ou práticas, onde as práticas no laboratório serão marcadas no decorrer da disciplina.}

\def\PUDavalia{A avaliação será feita com aplicação de provas teóricas e/ou práticas, além da possibilidade de inclusão de trabalhos e seminários no decorrer da disciplina.}

\def\PUDbibbasica{ &  \href{www.biblioteca.ifce.edu.br/index.asp?codigo_sophia=62499}{Carvalho}, M. M.; Paladini, E. P.  Gestão da qualidade: teoria e prática.  3ª ed.  São Paulo, SP: Atlas, 2012.  302p.  ISBN: 9788522471157. 
%\\ 
 %&  BANAS, F.;  Construindo Um Sitema de Gestão da Qualidade - Baseado na Norma ISO 9001-2008.  1ª edição.  São Paulo.  Editora EPSE.  2010.  312p.  ISBN 9788589705455.
\\
 &  \href{www.biblioteca.ifce.edu.br/index.asp?codigo_sophia=61854}{Carpinetti}, L. C. R.  Gestão da Qualidade - Conceitos e Técnicas.  3ª ed.  São Paulo, SP: Atlas, 2016.  247p.  ISBN: 9788597003918.
\\
 &  \href{www.biblioteca.ifce.edu.br/index.asp?codigo_sophia=65384}{Paoleschi}, B.  Logística Industrial Integrada - Do Planejamento, Produção, Custo e Qualidade a Satisfação do Cliente.  3ª ed.  São Paulo, SP: Érica, 2011.  264p.  ISBN: 9788536501970. }

\def\PUDbibcomplementar{ & \href{www.biblioteca.ifce.edu.br/index.asp?codigo_sophia=61899}{Carvalho}, M. M. de; Paladini, E. P.  Gestão da qualidade: teoria e casos.  2º ed.  Rio de Janeiro, RJ: Elsevier, 2012.  430p.  ISBN: 9788535248876.
%&  CARPINETTI, L.  C.  R.;  Gestão da Qualidade.  1ª edição.  São Paulo.  Editora ATLAS.  2010.  256p.  ISBN 9788522458028.
\\
 &  \href{www.biblioteca.ifce.edu.br/index.asp?codigo_sophia=24188}{Pereira}, Mário Jorge.  Engenharia de manutenção: teoria e prática.  Rio de Janeiro, RJ : Ciência Moderna, 2011. 228p.  ISBN: 9788573937879.
\\
 &  \href{www.biblioteca.ifce.edu.br/index.asp?codigo_sophia=24356}{Branco Filho}, Gil.  A Organização, o planejamento e o controle da manutenção.  Rio de Janeiro, RJ : Ciência Moderna, 2008. ISBN: 9788573936803.
\\
 &  \href{www.biblioteca.ifce.edu.br/index.asp?codigo_sophia=58562}{Seleme}, Robson; Stadler, Humberto.  Controle da qualidade: as ferramentas essenciais.  E-book.  Intersaberes.  186p.  ISBN: 9788565704861.
\\
 &  \href{www.biblioteca.ifce.edu.br/index.asp?codigo_sophia=40864}{Moura}, Luiz Antônio Abdalla de.  Qualidade e gestão ambiental: sustentabilidade e ISO 14.001.  6º ed.  Belo Horizonte, MG: Del Rey, 2014.  ISBN: 9788538401766. }

\def\PUDautor{Francisco Nélio Costa Freitas}

\def\PUDdata{2013-06-10}

\def\PUDrevisao{2}

\def\PUDrevisor{Weber Chaves Fontoura}

\def\PUDdatarevisao{2019-05-15}

\PUDcorpo
